\documentclass[11pt]{article}
\usepackage[T1]{fontenc}
\usepackage{microtype}
\usepackage[letterpaper,margin=1.08in]{geometry}
\usepackage{amsmath,amssymb,amsthm,mathtools}
\usepackage[hidelinks]{hyperref}
\newtheorem{theorem}{Theorem}[section]
\newtheorem{lemma}[theorem]{Lemma}
\newtheorem{corollary}[theorem]{Corollary}
\newtheorem{proposition}[theorem]{Proposition}
\newcommand{\CC}{\mathbb C}
\newcommand{\RR}{\mathbb R}
\newcommand{\QQ}{\mathbb Q}
\newcommand{\Gm}{\mathbb G_m}
\newcommand{\SL}{\mathrm{SL}}
\newcommand{\SU}{\mathrm{SU}}
\newcommand{\SO}{\mathrm{SO}}
\newcommand{\PGL}{\mathrm{PGL}}
\newcommand{\CT}{\operatorname{CT}}
\newcommand{\tr}{\operatorname{tr}}
\newcommand{\Lie}{\operatorname{Lie}}
\newcommand{\dd}{\,\mathrm d}
\DeclareMathOperator{\Ad}{Ad}
\title{Mathieu's Conjectures over Fields of Characteristic Zero\\
\large Split Homogeneous Spaces, Compact Groups, and Ground-Field Obstructions}
\author{Roy van Rijn}
\date{November 2026}
\begin{document}
\maketitle

\begin{abstract}
Let \(K\) be any field of characteristic zero. For an affine reductive
homogeneous space \(G/H\) with \(G^\circ\) split, we prove that its normalized
Reynolds kernel is Mathieu--Zhao exactly when \([G^\circ,G^\circ]\subseteq H\).
More generally, one nontrivial action of a split \(\SL_2\) gives a fixed
marked counterexample over \(K\), including when the acting group is
nonsplit. The proof transfers a rank-one witness along the principal
orbits of the action. A moment recurrence and Hilbert irreducibility
choose a ground-field parameter with infinitely many surviving marked
orders. We also classify the kernel on compact Hausdorff homogeneous
spaces and on connected reductive groups over real closed fields, and
give a ground-field nullcone criterion in the toral case.
\end{abstract}

\section{Introduction}

Let \(X=G/H\) be an affine homogeneous space of a linear algebraic group
over a field \(K\) of characteristic zero. Assume that \(G^\circ\) and
\(H^\circ\) are reductive. Its normalized Reynolds functional
\(\Lambda_X:K[X]\to K\) is the projection onto the invariant constants.
Call \(P\) \emph{pure} if \(\Lambda_X(P^m)=0\) for every \(m\ge1\).
It is \emph{safe} if every fixed \(Q\in K[X]\) satisfies
\(\Lambda_X(QP^m)=0\) for all sufficiently large \(m\).
The kernel is Mathieu--Zhao precisely when every pure function is safe.

Mathieu's first conjecture posed this implication for representative
functions on compact connected Lie groups \cite{Mathieu}. His second
asked whether pure functions on complex reductive spherical pairs are
unstable. Long classified the compact connected group case and the
complex spherical case \cite{Long,LongSpherical}. Over \(\CC\) the
classification of Reynolds kernels holds on arbitrary reductive
homogeneous spaces, with binary octahedral and icosahedral quotients
supplying the exceptional rank-one witnesses \cite{Complex}. The
ground-field question is different. The compact real quaternion
norm-one group has strictly positive squared trace on every nonzero
real regular function. Over the complex numbers its normalized trace
has a counterexample.

\begin{theorem}\label{thm:main}
Suppose that \(G^\circ\) is split over \(K\). Then
\[
 \ker\Lambda_{G/H}\text{ is Mathieu--Zhao}
 \quad\Longleftrightarrow\quad [G^\circ,G^\circ]\subseteq H .
\]
When containment fails there are \(P,Q\in K[G/H]\) such that \(P\) is
pure and \(\Lambda_{G/H}(QP^m)\ne0\) for infinitely many \(m\). For
connected \(G\) the same containment criterion characterizes equality
of the pure set with the geometric nullcone on \(K[G/H]\).
\end{theorem}

The negative direction is a consequence of an obstruction that depends
only on one subgroup. If a \(K\)-homomorphism from split \(\SL_2\) acts
nontrivially on \(X\), a general stabilizer either determines a
normalizer quotient of \(\SL_2\) or supplies an adjoint covariant.
A witness on the normalizer quotient can have poles after pullback.
Clearing them by an invariant denominator preserves purity. To preserve
marked nonvanishing over \(K\), we use a recurrence and choose the
denominator by arithmetic specialization.

For example, on the symmetric-matrix model of
\(\SL_3/\SO(2,1)\), the split subgroup \(\SO(2,1)\) acts by conjugation
after one fixed change of matrix coordinates. The general stabilizer
consists of sign changes of three distinct orthogonal eigenlines.
It fixes no adjoint vector. The unordered eigenlines nevertheless
give a rational map to the octahedral normalizer quotient, so the
normalizer argument applies.

We use the term \emph{geometric nullcone} for the points whose orbit
closure over an algebraic closure contains zero. For an individual
regular function this can be tested in any finite-dimensional rational
submodule containing it. For disconnected \(G\), geometric instability
is the same for \(G\) and \(G^\circ\), since the larger orbit closure is
a finite union of translates. An unstable function is safe. Over an
algebraic closure a contracting cocharacter gives strictly positive
weights on that function, eventually dominating the finitely many
weights of any fixed multiplier \cite{Kempf,LongSpherical}.
A nonzero marked subsequence therefore excludes geometric instability.

\section{Rank-one moment identities}\label{sec:rank}

Normalized Reynolds functionals commute with ground-field extension.
An equivariant morphism between transitive spaces preserves their
normalized functionals under pullback, by uniqueness. In particular,
on a quotient \(G/H\) the trace is the restriction of the normalized
left trace on \(K[G]\). We use these facts also for central isogenies
and products.

\begin{lemma}\label{lem:splitrank}
For split \(S=\SL_2\), every proper reductive stabilizer \(L\) gives
a pure regular function on \(S/L\) and one fixed multiplier with
infinitely many nonzero marked moments, over \(K\).
\end{lemma}

\begin{proof}
Use the binary forms
\[
 f_4=X^2Y^2,\qquad
 f_8=X^8+14X^4Y^4+Y^8,\qquad
 f_{12}=XY(X^{10}+11X^5Y^5-Y^{10}).
\]
For a form of degree \(n=2b\) write \(a=b-1\). For the rows \(u,v\)
of a determinant-one matrix let \(U=f(u)\), \(F=f(v)\), and set
\[
 R=\frac{f(v+Fu)}{F},\qquad P_f=U^aR^b,\qquad Q_f=F^a.
\]
Taylor expansion in the added row shows \(R\) is polynomial. Every
subgroup fixing \(f\) fixes the pair under right translation.
For \(f_4,f_8,f_{12}\), the groups are respectively the torus
normalizer, \(2O\) and \(2I\). The pure moments vanish over \(\QQ\).
The quartic has nonzero marked moments at all positive orders.
For the two exceptional forms there are nonzero marked moments at
orders \(m\) for which \(28m-3\), respectively \(66m-5\), is prime.
These are the formulas of \cite[Theorem~1.2 and Lemma~5.1]{Complex}.

The same conclusions hold for any nonzero scalar multiple of a conjugate
form of those degrees. Indeed, write \(A=am\), \(B=bm\), and
\(H_k^{(B)}=[t^k]f(v+tu)^B\). The left diagonal torus selects \(k=B\)
in the pure moment and \(k=B-b\) in the marked moment. The corresponding
marked summand is \(U^A F^{a+k-B}H_k^{(B)}\). Scaling \(f\) by \(\lambda\)
scales that summand by \(\lambda^{(n-1)m-1}\), and scales the one pure
summand by a nonzero scalar. Right translation by the conjugating
matrix preserves the left trace.

It remains to find such a form over \(K\). A nonzero vector in a
one-dimensional geometric invariant space can be taken over \(K\),
because taking invariants is a coaction kernel and commutes with base
change. For a positive-dimensional proper \(L\), use the quartic of
the normalizer of \(L^\circ\), which is a torus. For a noncentral
finite cyclic group use the centralizer torus. For a binary dihedral
group of order \(4r\) with \(r\ge3\), the cyclic subgroup generated by
elements of order greater than four is characteristic. Its centralizer
is a torus normalized by \(L\). For a central \(L\) choose any split
torus. In each case the torus normalizer has a unique invariant line
of degree four. For \(Q_8\), use its normalizer, which over the algebraic
closure is \(2O\). Indeed \(Q_8\) is normal in \(2O\), its centralizer
is the center, and its automorphism group has order \(24\).
For \(2T,2O,2I\), use degrees \(8,8,12\), respectively.
The octic is invariant under \(2T\subset 2O\). Uniqueness of these
lines follows because the ratio of two independent invariants of
degree \(n\) would give a nonconstant map \(\mathbb P^1\to\mathbb P^1\)
of degree at most \(n\) with a full generic projective group orbit in
each generic fiber. The projective orders \(12,24,60\) exceed the
respective degrees. This exhausts the finite cases by the classification
in \cite{Flicker}.
\end{proof}

For use when a principal stabilizer is central, write a traceless
matrix as
\[
 Z=\begin{pmatrix}a/2&v\\u&-a/2\end{pmatrix},\quad r=a^2+4uv,
 \qquad P_{\rm ad}=v^2(1+au-u^3v)^4,\quad Q_{\rm ad}=u^2 .
\]
If \(R_S\) denotes the conjugation Reynolds projection, then
\begin{equation}\label{eq:adjoint}
 R_S(P_{\rm ad}^m)=0,\qquad
 R_S(Q_{\rm ad}P_{\rm ad}^m)=\beta_m r^{3m-1},\qquad
 \beta_m=\frac{2m}{3m+1}\frac{(2m)!(4m)!}{(6m)!}.
\end{equation}
To check this, take \(Z=ww^*-\|w\|^2I/2\) with \(w=(x,y)^{\mathsf T}\).
Then \(u=\bar x y\), \(v=x\bar y\), \(a=|x|^2-|y|^2\).
The quartic normalizer pair on the matrix with columns
\(w,(-\bar y,\bar x)^{\mathsf T}\) becomes the displayed adjoint pair,
because its two parenthesized factors multiply to
\((1+|x|^2u)(1-|y|^2u)=1+au-u^3v\).
The radial identity of \cite[Section~5]{Complex} gives
\eqref{eq:adjoint}, using \(r(Z)=\|w\|^4\). The Hermitian matrices
obtained this way are all the traceless Hermitian matrices, a
Zariski-dense real form. Thus the identities of polynomial Reynolds
projections hold over every characteristic-zero field.

\section{Principal stabilizers}\label{sec:principal}

Let a split \(S=\SL_2\) act nontrivially on \(X\) through \(G^\circ\),
with \(G\) connected for the moment. We use two facts about the action.
There is a dense invariant open \(U\) with closed geometric orbits
and a single geometric conjugacy class of stabilizers. If that class
is noncentral, there is a rational equivariant map \(X\dashrightarrow
S/N_S(L)\), where \(L\) is the stabilizer of some \(K\)-point of \(U\).

For the first fact it is enough to use a finitely generated model
embedded in \(\CC\). Choose a compact form \(C\) of \(G_\CC\) containing
the standard compact subgroup of \(S_\CC\). A compact form of \(H_\CC\)
can be conjugated into \(C\), and we may replace \(H_\CC\) by that
conjugate using a \(G_\CC\)-equivariant isomorphism of quotients
\cite[Theorems~3.7--3.8]{AdamsTaibi}. The compact conjugation then
defines an antiholomorphic involution on \(X_\CC\), with a dense
compact orbit of fixed points. Choose an embedding into a
finite-dimensional \(S_\CC\)-module compatible with that
conjugation, by taking an invariant space of algebra generators
closed under it. In a compact invariant Hermitian metric real on
the fixed real form, every real vector has zero norm derivative
in the \(i\Lie(\SU_2)\) directions. Its orbit is closed by
Kempf--Ness \cite{KempfNess}. The compact orbit therefore meets both
the principal Luna stratum and the maximal-dimension orbit locus
\cite{Luna}. On the principal stratum a nonclosed orbit would
have to be larger than its closed orbit, whose dimension is
already maximal. Thus the inverse image of this stratum is the
desired open. Uniqueness and base change of invariant algebras
give descent to \(K\).

For the second fact, \(G(K)\) is Zariski dense. One way to see this is
to choose regular semisimple Lie algebra vectors spanning \(\Lie(G)\)
and take the associated maximal \(K\)-tori. Each torus has dense
\(K\)-points by the norm map from restriction of scalars of a split
torus. Multiplication of the tori is dominant by its differential at
the identity. Thus \(U\) has a \(K\)-point \(x_0\). Set \(L=S_{x_0}\).
It is reductive by the affine Matsushima criterion \cite{Matsushima}.
If \(L\) is noncentral, \(N=N_S(L)\) is proper and reductive, as the
rank-one subgroup list shows. For a finite noncyclic \(L\) its
centralizer and normalizer are finite. For a noncentral cyclic \(L\),
the normalizer preserves the centralizer torus. If \(L^\circ\) is
a torus, the normalizer preserves that torus. The map in question
sends a point with stabilizer \(sLs^{-1}\) to \(sN\). More formally,
the image of
\[
 S\times U^L\longrightarrow U\times S/N,\qquad
 (s,y)\longmapsto(sy,sN)
\]
is a constructible graph with one geometric point over each point
of \(U\). The dominating component of its closure is birational
to \(U\) by generic separability in characteristic zero. This
gives the rational map \(\phi\) over \(K\).

\subsection*{The two types of witness}

If \(L\) is central, restriction to a closed orbit shows that the
adjoint module occurs in \(K[X]\). Restriction is surjective over
the algebraic closure, the adjoint module occurs on \(S/L\), and
the intertwiner space commutes with base change. Hence there is
a \(K\)-defined nonzero equivariant map to the adjoint space.
There is also such a map over \(\CC\) taking Hermitian values on the
compact orbit. In fact the antilinear involution
\(
 Z\mapsto \tau_{\rm h}\circ Z\circ\tau_X
\)
on the space of equivariant maps uses the compact conjugation
\(\tau_X\) of the quotient and \(\tau_{\rm h}(M)=\bar M^{\mathsf T}\).
Its real fixed space spans the complex space. Choose a nonzero
fixed map. It has nonzero radius on an open subset of the compact
orbit. By \eqref{eq:adjoint}, therefore, all its positive marked
moments are nonzero. It can be written in a finite span of maps
defined over \(K\). This gives a finite-parameter family of pure
adjoint phases over \(K\), whose marked sequences are generically
nonzero.

If \(L\) is noncentral, pull back the witnesses \(P_0,Q_0\) of
Lemma~\ref{lem:splitrank} by \(\phi\). Denote the two rational functions
by \(p,q\). Their translates span a finite-dimensional rational
submodule. Its common denominator ideal \(I\) in \(K[X]\) is nonzero
and \(S\)-stable. Over \(\CC\), choose a nonzero \(a\in I\). Write
\(\tau_X^*a\) for compact conjugation on functions.
Then \(d=R_S(a\,\tau_X^*a)\) belongs to \(I^S\).
On the compact orbit it is nonnegative, and it is positive somewhere
by density and the compact angular average. Hence
\(\Lambda_X(d^j)\ne0\) for every \(j\ge1\). The denominator ideal
and its invariant subspace commute with field extension. Indeed,
after choosing a finite basis \(a_j/b\) with a common denominator
\(b\), the ideal is the simultaneous kernel of multiplication
by the \(a_j\) into \(K[X]/(b)\). Therefore we may take a finite
parameter span of \(K\)-invariant denominators containing \(d\)
after field extension. For each member of that span the regular pair
\((dp,dq)\) is pure and has marked moments
\begin{equation}\label{eq:factor}
 \Lambda_X((dq)(dp)^m)
    =\Lambda_X(d^{m+1})\,\Lambda_{S/N}(Q_0P_0^m).
\end{equation}
Indeed, on the open of definition the \(S\)-projection gives constant
rank-one moments, times \(d^{m+1}\) in the marked case. The identities
extend by regularity, and \(\Lambda_X\circ R_S=\Lambda_X\).
The nonzero complex parameter above shows that the polynomial sequence
in \eqref{eq:factor} is not eventually zero.

\section{Arithmetic specialization}\label{sec:arithmetic}

\begin{lemma}\label{lem:specialize}
Let \(E/\QQ\) be finitely generated. Suppose
\(b_m(\boldsymbol u)\in E[\boldsymbol u]\) is not eventually zero over
\(E(\boldsymbol u)\) and satisfies a scalar polynomial recurrence
with nonzero first coefficient for all sufficiently large integer
indices. Then some specialization \(\boldsymbol u\in E^d\) keeps
infinitely many \(b_m\) nonzero.
\end{lemma}

\begin{proof}
Write \(A(m,\boldsymbol u)\) for the first coefficient. Remove its
content in \(E[m]\) when tracking its zeros, and choose \(M\) beyond
the fixed integer zeros and the recurrence cutoff. Choose
\(m_0\ge M\) with \(b_{m_0}\ne0\) as a polynomial. On a line through
two suitably chosen \(E\)-points the reduced coefficient has no
constant root in \(m\), and \(b_{m_0}\) is still nonzero. To see this,
choose the first point with \(b_{m_0}\ne0\) and reduced coefficient
a nonzero polynomial in \(m\). Choose the second point to avoid
vanishing of the reduced coefficient at the finitely many roots
of the first evaluation. Each of these avoidance conditions is
possible since the content was removed. Denote the line parameter
by \(u\).

Factor the reduced coefficient in \(E(u)[m]\). Each linear factor
has a nonconstant root \(r_i(u)\). Choose a distinct rational
prime \(p_i\) and embed \(E\) into \(\QQ_{p_i}\) so that \(r_i(u)\)
has negative valuation on a nonempty \(p_i\)-adic open set. For
a finite pole \(\alpha\) of \(r_i\), embed \(E(\alpha)\) instead,
then take a small ball near the pole. A pole at infinity requires
only large \(p_i\)-adic absolute value. The embeddings exist for
infinitely many primes for any finitely generated field of
characteristic zero. Indeed, specialize a transcendence basis to
rational values where a minimal polynomial of a primitive element
has no coefficient poles or multiple roots. A root of the specialized
polynomial is algebraic over \(\QQ\), so it has a simple root in
\(\QQ_p\) for infinitely many \(p\), by Schur's prime-divisor theorem
and Hensel's lemma. Perturb the rational parameters to algebraically
independent values sufficiently close in \(\QQ_p\), and lift the
simple root there to get the required embedding.

Choose a rational \(a\) in all the required local opens by weak
approximation. Put \(b=\prod_i p_i\), take rational \(c\ne0\)
sufficiently small at these primes, and choose a prime \(\ell\)
larger than the degree of \(A\) in \(m\). Make the substitution
\[
                         u=a+\frac{c}{1+b t^\ell}.
\]
For every \(t\in E\), the reciprocal denominator is integral at the
chosen places: the summands have unequal valuations, since
\(v_{p_i}(b)=1\). Thus the local conditions always hold. The extension
\(E(t)/E(u)\) has degree \(\ell\). By coprimality it preserves
irreducibility of every factor in \(m\) of degree at least two.
If no local conditions are needed, take \(u=t\). The field \(E\) is
Hilbertian \cite[Chapter~13]{FriedJarden}. Specialize \(t\) to
preserve the degrees and irreducibility of the nonlinear factors,
avoiding the finitely many poles and zeros of \(b_{m_0}\). Linear
roots have negative valuation at one assigned prime and hence
cannot be integers. The resulting first coefficient is nonzero
at every integer \(m\ge M\). Backward recurrence from an eventual
zero tail would give \(b_{m_0}=0\), a contradiction.
\end{proof}

\begin{theorem}\label{thm:active}
Over every characteristic-zero field, a nontrivial action of a split
\(\SL_2\) on an affine reductive homogeneous space of a connected group
gives a fixed marked counterexample to the Mathieu--Zhao property.
\end{theorem}

\begin{proof}
Take a finitely generated field of definition for the preceding
constructions and finite parameter spans. In the central case it
suffices to specialize the sequences
\(\Lambda_X(r(Z)^{3m-1})\), since the factors \(\beta_m\)
in \eqref{eq:adjoint} are nonzero. In the noncentral case use the
whole product sequence \eqref{eq:factor}, so that the possibility
of gaps in the rank-one marked sequence is included. Each sequence
is polynomial in the parameters and is generically not eventually
zero.

These scalar sequences satisfy polynomial recurrences over the
parameter field. They are mixed moments of a regular function on
the acting group, or on a product of acting groups, by pullback.
On a connected group the invariant volume form \(\omega\) identifies
the trace with a functional on top forms annihilating exact forms.
Infinitesimal left translations give a frame preserving \(\omega\)
and integrate to zero. For a nonzero phase \(F\) over the parameter
field, use top cohomology of the connection
\(\mathrm d+\nu\,\mathrm dF/F\wedge\) on \(F\ne0\).
It is finite-dimensional over the rational field in \(\nu\) by
the algebraic Mellin shift \cite[Section~1.3]{LoeserSabbah}.
The classes of \(B\omega,FB\omega,\ldots\), for a fixed multiplier \(B\),
are therefore linearly dependent. At a sufficiently large integer
\(m\), multiplication by \(F^m\) turns each twisted exact term into
an ordinary regular exact form. This gives the recurrence for
\(\Lambda(BF^m)\). After clearing denominators the relation is
polynomial in the parameters for all sufficiently large integers.
In \eqref{eq:factor} use phase \(d\otimes P_0\) and multiplier
\(d\otimes Q_0\). In the central case use \(r^3,r^2\) and shift
the index by one. Lemma~\ref{lem:specialize} now gives the asserted
fixed pair over the finitely generated field and hence over \(K\).
\end{proof}

\begin{proof}[Proof of Theorem~\ref{thm:main}]
Suppose \(D=[G^\circ,G^\circ]\) is not contained in \(H\). On the
\(K\)-defined open-and-closed component \(G^\circ H/H\) some split
almost simple factor acts nontrivially, after a central cover.
A root subgroup of an active factor gives an action as in
Theorem~\ref{thm:active}. Extend the pair by zero to the other
components. Its marked moments are multiplied by the reciprocal
of the number of geometric components. If \(D\subseteq H\), the
geometric components are tori. The component average and
finite-trace separation of \cite[Section~7]{Complex}, followed
by the Laurent theorem \cite{DvK}, prove the Mathieu--Zhao property
over the algebraic closure, hence over \(K\). For connected \(G\)
the Laurent support half-space criterion also proves geometric
instability of pure functions in this toral case. The opposite
case has a marked witness outside the nullcone by the weight bound.
\end{proof}

\section{Real closed fields and the nonsplit boundary}\label{sec:real}

\begin{theorem}\label{thm:realgroup}
Let \(K\) be real closed and \(G\) connected reductive. The normalized
trace on \(K[G]\) has Mathieu--Zhao kernel exactly when \(G\) is a torus
or \(G\) is \(K\)-anisotropic.
\end{theorem}

For the positive compact case, a nonzero \(f\) has nonzero squared
trace. One can embed a finitely generated field of definition of the
group and \(f\) into \(\RR\), preserving the negative definiteness of
the derived Killing form and the anisotropic central torus. For the
first condition preserve the signs of the principal minors. The
anisotropic central torus over a real closed field is a product of
norm-one tori for \(K(i)/K\), so include this model isomorphism in
the field of definition. On the compact real identity group
\(\Lambda(f^2)>0\) by density of the real orbit and normalized Haar.
This also proves the positive assertion on every affine \(G/H\) with
such a connected compact acting group. The torus case follows from
\cite{DvK} after splitting.

We record the real specialization principle used here and below.
Given finitely many strict signs on elements of a finitely generated
ordered field, embed that field into \(\RR\) respecting those signs.
Write the field as a finite separable extension of
\(\QQ(t_1,\ldots,t_d)\). Express existence of the root, the signs,
nonzero denominators and the nonzero derivative by a finite formula.
It holds in a real closure, hence over \(\RR\) by quantifier
elimination \cite{Tarski}. The simple root persists on a real open
set of parameters. Choosing algebraically independent values in
that open gives the desired embedding. If \(d=0\), the simple real
root itself suffices. In particular a finite condition with
parameters in a real closed field follows if it holds at every real
embedding of a finitely generated field of definition. Otherwise
quantifier elimination would give finitely many signs certifying a
failure, which could be preserved by such an embedding.

For the first negative case, a noncompact derived group has a split
rank-one subgroup, so Theorem~\ref{thm:active} applies to \(G\).
The same argument applies to any \(G/H\) on which a noncompact almost
simple factor acts nontrivially. Indeed a nontrivial split
cocharacter in the derived group gives a rational unipotent and
hence a split rank-one homomorphism by rational Jacobson--Morozov
\cite{JM}. Over a real closed field the noncompact derived case can
equivalently be recognized by a positive derived Killing square.
Over \(\RR\), a vector in the noncompact summand of a Cartan
decomposition has self-adjoint adjoint endomorphism. A nonzero
eigenvector of nonzero eigenvalue is adjoint-nilpotent by weight
shift. The existence of a nonzero \(\mathfrak{sl}_2\)-triple is a
finite condition, so the real specialization principle gives
the same conclusion over \(K\).

The other negative case of Theorem~\ref{thm:realgroup} has nontrivial
compact derived group and split central direction. Here is a
ground-field construction that also works over arbitrary fields.

\begin{proposition}\label{prop:central}
Suppose \(G\) is connected reductive, it has a nontrivial \(K\)-character
\(\chi\), and its derived group contains a \(K\)-subgroup \(S\) of
geometric type \(A_1\). Then there is a fixed marked counterexample
on \(K[G]\).
\end{proposition}

\begin{proof}
Write \(q\) for the negative of the derived Killing form on
\(\mathfrak s=\Lie(S)\), with polarization \(\langle\,,\,\rangle_q\).
Its restriction is nondegenerate, as one sees after splitting \(S\)
by evaluating the sum of squared weights on a coroot. Let \(p\) be
the Killing-orthogonal projection to \(\mathfrak s\), choose
\(\xi\in\mathfrak s(K)\) with \(q(\xi)\ne0\), and set
\[
 z=p(\Ad(g)\xi),\quad A=q(z),\quad
 B=\frac{\langle\xi,z\rangle_q^2}{q(\xi)},\quad t=\chi(g).
\]
Use
\begin{equation}\label{eq:angular}
 P=t^{-1}(1+t)\{A-(1+t)^2B\}^{2},\qquad Q=t .
\end{equation}
To check moments, use a model embedded in \(\CC\). By conjugating
the compact form of \(S\) we may write \(\xi=c\xi_0\) with \(\xi_0\)
compact and \(c\ne0\). This uses transitivity of the adjoint action
on each nonzero quadric level. Extend that compact form to the
derived group. Replacing \(\xi_0\) by \(c\xi_0\) scales \(P\) by
\(c^4\). For \(\xi_0\), the law of \(z\) under compact derived Haar
measure is rotation-invariant in the three-dimensional compact
adjoint space. On each nonzero radius \(B/A=w^2\) with \(w\)
uniform on \([0,1]\). The central isogeny from the product of the
connected center and the derived group separates the trace.
The variable \(t\) is a nontrivial central character, so central
averaging extracts its constant term.

Put \(F_m(s)=\int_0^s(1-w^2)^{2m}\dd w\). The angular average of
\(P^m/A^{2m}\) is
\[
                  t^{-m}(1+t)^{m-1}F_m(1+t).
\]
Since \(F_m(1+t)-F_m(1)\) is divisible by \(t^{2m+1}\), its constant
term is zero. Its constant term after multiplication by \(t\)
is \(F_m(1)>0\). The remaining radial integral of \(A^{2m}\)
is positive, since \(z\) is nonzero at the identity. This proves
the two moment statements after rescaling and hence over \(K\).
The same pair descends to \(G/H\) if \(H\) fixes \(\xi\) and
\(\chi|_H=1\).
\end{proof}

For a nontrivial compact derived group over a real closed field,
a maximal torus in it is anisotropic. Conjugation acts by \(-1\)
on its character lattice, so an opposite pair of roots generates
a rank-one subgroup over \(K\). If the center has a split direction,
there is a nontrivial \(K\)-character, and Proposition~\ref{prop:central}
applies. This completes Theorem~\ref{thm:realgroup}.

\subsection*{The two compact products}

For connected homogeneous spaces over real closed fields,
Theorem~\ref{thm:active} handles the active noncompact factors.
Remove inactive factors. When the remaining derived group is
compact, a central cover expresses the group as a split torus
times a compact-type group. The projection of the connected
stabilizer to the split factor is contained in the stabilizer:
the split part of its connected center maps trivially to the
compact group, and its anisotropic part has no split characters.
Quotient out that split subtorus. In the resulting product
\(T\times C\), the stabilizer \(H'\) has finite projection
\(\Gamma\) to \(T\). Write \(L'\) for its projection to \(C\).
There is a homogeneous quotient map to
\((T/\Gamma)\times(C/L')\). The compact derived group is active
there whenever it was active before, because \(H'\cap C\) has
finite index in \(L'\).

If one of the split and derived directions is absent, square
positivity or the Laurent theorem applies. Otherwise use a
character on \(T/\Gamma\) and an active compact rank-one subgroup
of \(C\). The principal-orbit argument applies to \(C/L'\).
For a central principal stabilizer there is an adjoint covariant
with values in the compact Lie algebra and the angular construction
\eqref{eq:angular} applies. For a noncentral stabilizer there is
a rational map to its normalizer quotient. If that normalizer
lies in a torus normalizer, pull back the axis-square angular
witness on the unoriented two-sphere. An invariant denominator
clears the finitely many Laurent coefficients and preserves the
marked moments by compact square positivity. The compact
rank-one normalizer list leaves two products,
\[
                 (\SU_2/2O)\times\Gm,\qquad
                 (\SU_2/2I)\times\Gm ,
\]
where \(\SU_2\) is the compact real algebraic form. A finite
stabilizer at a point of the compact orbit lies in the compact
rank-one group by the norm-minimality argument and polar decomposition.
Its noncyclic finite normalizer is likewise compact, since the
positive part of a normalizing element centralizes the finite
subgroup. Thus the products use the usual compact finite subgroups.

For these two products the left circle quotient pushes Haar to
uniform area on the two-sphere modulo the projective finite rotation
group. Adjoining inversion gives the reflection group \(B_3\) or
\(H_3\). In its two nonradial basic invariant coordinates the density
on the regular image is proportional to the inverse square root of
the discriminant. Indeed the tangential Jacobian differs from the
full Jacobian of the three basic invariants by the constant
length of the gradient of the radial quadratic on the sphere.
The full Jacobian is the product of reflecting linear forms
\cite{Chevalley,Steinberg}. This identifies the density that a
real Laurent normalizer on these coordinates must handle. The direct
substitution \(p=t^{-1}(1+t)f((1+t)x)^2\), with \(f\) the discriminant,
already has a nonzero first moment for this density. Indeed integrate
the divergence of \(x f^{3/2}\) on the bounded two-dimensional
image. The field vanishes on the boundary and has bounded first
derivatives, so its divergence integrates to zero. Hence
\[
 \int \bigl(f^2+2f\,x\cdot\nabla f\bigr)f^{-1/2}\dd x
   =-\frac53\int f^{3/2}\dd x .
\]
The left side is the unnormalized first constant-term moment of \(p\).
This accounts for the failure of the unmodified polynomial-density
seed.

\subsection*{An anisotropic arithmetic test case}

The split rank-one obstruction does not decide anisotropic
individual forms over a general field. For example, the quaternion
algebra \((2,3)_{\QQ}\) is division by its Hilbert symbol at \(3\),
but it splits over \(\RR\). Its norm-one group is anisotropic over
\(\QQ\), and its real trace is outside the compact positivity
case. On any anisotropic norm-one quadric, two simple witness
classes can be ruled out. A nonzero homogeneous linear phase
has nonzero second moment by the inverse norm form. For the even
degree-two class, absorb the constant into the norm equation and
write the phase as \(q_A(x)=q(x,Ax)\), where \(A\) is self-adjoint
for \(q\). In dimension \(n\), its \(m\)-th normalized sphere
moment equals
\[
 \frac{m!}{(n/2)_m}[z^m]\det(I-zA)^{-1/2}.
\]
This follows by diagonalization on a dense open of the complex
self-adjoint matrices and the Gaussian radial integral, then by
polynomial continuation. If all positive moments vanish, \(A\)
is nilpotent. A nonzero self-adjoint nilpotent would give a
nonzero isotropic vector: take \(A^r v\ne0\) with \(2r\) at least
the nilpotence index. Thus \(A=0\). This calculation leaves
inhomogeneous and higher-degree phases open.

\section{Compact homogeneous spaces}

\begin{theorem}\label{thm:compact}
Let \(C\) be a compact Hausdorff group and \(L\) a closed subgroup.
On the complex representative functions of \(C/L\), normalized
invariant integration has Mathieu--Zhao kernel exactly when
\(\overline{[C^\circ,C^\circ]}\subseteq L\). In particular on \(C\)
it holds exactly when \(C^\circ\) is abelian.
\end{theorem}

\begin{proof}
In the compact Lie case complexify the pair. A faithful unitary
model has reductive Zariski closure, and its regular functions
restrict to all representative functions. The matrix-entry algebra
is conjugation-stable on a unitary group, so it is dense by
Stone--Weierstrass. Projecting onto each finite-dimensional
Peter--Weyl isotype proves the asserted equality. The Zariski
closure of \(L\) meets the compact Lie group in \(L\): average
matrix-entry polynomials over \(L\) to separate distinct right
cosets. Right invariance and normalized integration therefore agree
with the complex homogeneous-space constructions. The identity
component of a compact Lie group is Zariski dense in the corresponding
algebraic identity component, and the same holds for its derived
subgroup. Apply the complex classification.

For a general compact Hausdorff group, any two representative functions
factor through a single compact Lie quotient \(J\), using the direct
sum of their finite-dimensional representations. If the functions
are right \(L\)-invariant, they descend to \(J/L_J\) with
\(L_J\) the image of \(L\). The image of \(C^\circ\) is \(J^\circ\)
because the quotient by that image is a quotient of the profinite
group \(C/C^\circ\). Haar measure pushes forward, so the containment
criterion suffices by the compact Lie case. Conversely, an element
of \([C^\circ,C^\circ]\) outside \(L\) can be separated from \(L\)
by a finite sum of representations, using point separation and
compactness of \(L\). The resulting compact Lie quotient provides
a counterexample which pulls back \cite{Fausk}.
\end{proof}

If one defines instability using the Zariski closure of the compact
orbit in the finite-dimensional complex orbit span, the corresponding
nullcone criterion holds precisely when the lattice criterion of
\cite[Section~7]{Complex} holds on every compact Lie quotient pair.
Both the moments and the orbit of an individual representative
function factor through such a quotient. In particular for connected
\(C\), the same containment criterion governs both properties.

\section{Universal Reynolds projections and toral nullcones}

\begin{theorem}\label{thm:universal}
For a linear algebraic group \(G\) with reductive identity component
over \(K\), the following are equivalent.
\textup{(i)} \(G^\circ\) is a torus.
\textup{(ii)} The Reynolds projection has Mathieu--Zhao kernel on every
commutative rational \(G\)-algebra.
\textup{(iii)} For every such algebra and every \(f\) with zero Reynolds
projection of each positive power, each fixed geometric isotypic
component of \(f^m\) vanishes eventually.
\end{theorem}

\begin{proof}
Assume \(G^\circ=T\) is a torus. Work over the algebraic closure
in the finitely generated stable subalgebra \(B\) generated by a
finite-dimensional submodule containing \(f\). Choose component
representatives \(g_1=1,\ldots,g_s\), and write \(L_i\) for the
Laurent polynomial recording the \(T\)-weights of \(g_i f\) with
coefficients in \(B\). Since the component quotient acts by averaging
the \(T\)-projection, \(\sum_i\CT L_i^m=0\) for every \(m\ge1\).

Take a primary decomposition \(0=\bigcap_j\mathfrak q_j\) in \(B\).
Put \(\mathfrak p_j=\sqrt{\mathfrak q_j}\), and choose \(e_j\) with
\(\mathfrak p_j^{e_j}\subseteq\mathfrak q_j\).
In the fraction field of \(B/\mathfrak p_j\), finite-trace separation
\cite[Lemma~7.1]{Complex} gives pure constant terms for each Laurent
polynomial separately. Write \(L_1=U+V\) by putting the coefficients
in \(\mathfrak p_j\) in \(V\). If \(U\ne0\), the Laurent theorem gives
a functional strictly positive on its support. Terms of \((U+V)^m\)
with fewer than \(e_j\) factors of \(V\) then eventually cannot
reach a prescribed weight. All other terms have coefficients in
\(\mathfrak q_j\). The case \(U=0\) is immediate. Thus each prescribed
weight of \(f^m\) is eventually zero by the intersection of the
primary components. This proves both conclusions: each fixed
multiplier and each fixed geometric simple module has only finitely
many \(T\)-weights.

If \(G^\circ\) is non-toral, Theorem~\ref{thm:main} after algebraic
closure gives a pair on the group that is defined over a finite
extension \(E/K\). Regard \(E[G_E]\) as a rational \(G\)-algebra
over \(K\). Its invariants are \(E\), and its Reynolds projection
is the normalized trace with values in \(E\). The pair disproves
the universal kernel assertion. It also disproves eventual isotypic
vanishing, since the trace of a product pairs only dual types and
a fixed multiplier has finitely many types.
\end{proof}

Finally suppose \(D=[G^\circ,G^\circ]\subseteq H\), without a splitness
assumption. Let \(T=G^\circ/D\), \(M=X^*(T_{\overline K})\), and let
\(\Gamma\) be a finite Galois group acting on \(M\) and the geometric
components. For a representative \(i\) of each component orbit
\(\alpha\), let \(\Gamma_i\) be its stabilizer and \(M_i\subset M\)
the weight lattice of that component. Put
\(V_i=(M_i\otimes\QQ)^{\Gamma_i}\), \(V=(M\otimes\QQ)^\Gamma\).
The ground-field pure set equals the geometric nullcone on \(K[X]\)
exactly when
\begin{equation}\label{eq:transfer}
 \bigoplus_\alpha V_i\longrightarrow V,\qquad
 (v_i)\longmapsto\sum_\alpha|\Gamma|^{-1}
                                  \sum_{\gamma\in\Gamma}\gamma v_i
\end{equation}
is injective. In fact purity separates componentwise over
\(\overline K\). The convex hull of each component support avoids
zero, and this remains true for its stabilizer-averaged weights
because the support is \(\Gamma_i\)-stable. If \eqref{eq:transfer}
is injective, strictly positive functionals on these averaged
supports combine into one invariant functional on \(M\otimes\QQ\).
It is positive on the original weights and yields a contracting
cocharacter after clearing denominators and lifting through the
connected center. Conversely, a nonzero integral relation in
\eqref{eq:transfer} gives one nonzero weight on each of the
participating Galois orbits of components. The corresponding
direct sum of character lines descends, and a nonzero \(K\)-vector
on each participating orbit is nonzero on every geometric
component of that orbit. This produces a pure function whose
global support has zero in its convex hull. It is therefore
outside the geometric nullcone. For split \(T\) this uses one
embedded character lattice per Galois orbit.

\subsection*{Two ground-field transfers}

Let \(E/K\) be finite separable. A fixed pair on \(Y=G/H\) over \(E\)
gives a fixed pair on its Weil restriction to \(K\). After scalar
extension to \(\overline K\), that restriction is the product of
the conjugates \(Y^\sigma\). The products of the conjugate phases
and multipliers are defined over \(K\), and their marked moment
is \(N_{E/K}(\Lambda_Y(QP^m))\). Purity and the precise nonzero
marked orders therefore descend to this restricted space.

Disconnected forms can introduce an additional field issue. For
a real closed \(K\), let \(C\) be the compact quaternion norm-one
group and take \(G=C\times\mu_3\). The two nonreal components are
represented by the \(K(i)[C_{K(i)}]\) factor of \(K[G]\).
Use on it a complex rank-one pair with rational marked moments
and extend by zero on the other factor. The component average
gives a bad pair over \(K\) with marked moments scaled by \(2/3\).
Thus the connected group classification over real closed fields
does not extend by inspecting the real identity component alone.

\paragraph{AI assistance and review.}
ChatGPT (OpenAI) assisted with mathematical arguments, literature checks,
exposition, proof review and verification code.

\begin{thebibliography}{99}
\small
\setlength{\itemsep}{2pt}
\setlength{\parsep}{0pt}
\bibitem{Mathieu}
O. Mathieu, Some conjectures about invariant theory and their applications,
\emph{S\'emin. Congr.} 2 (1997), 263--279.
\href{https://smf.emath.fr/sites/default/files/2018-11/smf_sem-cong_2_263-279__sample.pdf}{Publisher excerpt}.

\bibitem{Long}
C. D. Long, The Mathieu property for compact connected Lie groups,
\href{https://arxiv.org/abs/2609.16178}{arXiv:2609.16178}.

\bibitem{LongSpherical}
C. D. Long, A classification of Mathieu's second conjecture for
reductive spherical pairs, September 2026,
\href{https://github.com/long-mathematics/mathieu-second-conjecture-spherical-pairs/blob/e9edac262b4b64eb68feb0e6e80496cefe4a1dee/mathieu_second_conjecture_spherical_pairs.tex}{versioned source e9edac2}.

\bibitem{Complex}
R. van Rijn, Counterexamples to the Mathieu property: from explicit
quotient constructions to a general classification, 2026.
\href{https://hexagonmath.org/2611.00556v1}{Hexagon 2611.00556v1},
\href{https://doi.org/10.5281/zenodo.24310836}{doi:10.5281/zenodo.24310836}.

\bibitem{DvK}
J. J. Duistermaat and W. van der Kallen, Constant terms in powers of
a Laurent polynomial, \emph{Indag. Math.} 9 (1998), 221--231.
\href{https://doi.org/10.1016/S0019-3577(98)80020-7}{DOI}.

\bibitem{Flicker}
Y. Z. Flicker, Finite subgroups of \(\SL(2,F)\) and automorphy,
\emph{Manuscripta Math.} 167 (2022), 213--243.
\href{https://doi.org/10.1007/s00229-021-01274-x}{DOI}.

\bibitem{Luna}
D. Luna, Slices \'etales, \emph{M\'em. Soc. Math. France} 33 (1973),
81--105. \href{https://numdam.org/articles/10.24033/msmf.110/}{Numdam}.

\bibitem{KempfNess}
G. Kempf and L. Ness, The length of vectors in representation spaces,
\emph{Lecture Notes in Math.} 732 (1979), 233--243.
\href{https://doi.org/10.1007/BFb0066647}{DOI}.

\bibitem{Kempf}
G. Kempf, Instability in invariant theory, \emph{Ann. of Math.}
108 (1978), 299--316. \href{https://arxiv.org/abs/1807.02890}{Transcription}.

\bibitem{AdamsTaibi}
J. Adams and O. Ta\"ibi, Galois and Cartan cohomology of real groups,
\emph{Duke Math. J.} 167 (2018), 1059--1097.
\href{https://arxiv.org/abs/1611.05956}{arXiv:1611.05956}.

\bibitem{Matsushima}
I. V. Arzhantsev, Invariant ideals and Matsushima's criterion,
\href{https://arxiv.org/abs/math/0506430}{arXiv:math/0506430}.

\bibitem{FriedJarden}
M. D. Fried and M. Jarden, \emph{Field Arithmetic}, fourth ed.,
Springer, 2023. \href{https://doi.org/10.1007/978-3-031-28020-7}{DOI}.

\bibitem{LoeserSabbah}
F. Loeser and C. Sabbah, \'Equations aux diff\'erences finies et
d\'eterminants d'int\'egrales de fonctions multiformes,
\emph{Comment. Math. Helv.} 66 (1991), 458--503.
\href{https://perso.pages.math.cnrs.fr/users/claude.sabbah/articles/loeser-sabbah91.pdf}{PDF}.

\bibitem{Tarski}
C. Cohen and A. Mahboubi, Formal proofs in real algebraic geometry:
from ordered fields to quantifier elimination,
\emph{Log. Methods Comput. Sci.} 8 (2012).
\href{https://arxiv.org/abs/1201.3731}{arXiv:1201.3731}.

\bibitem{JM}
R. Appenzeller, Semisimple algebraic groups over real closed fields,
\href{https://arxiv.org/abs/2601.07732}{arXiv:2601.07732}, \S5.8.

\bibitem{Chevalley}
C. Chevalley, Invariants of finite groups generated by reflections,
\emph{Amer. J. Math.} 77 (1955), 778--782.

\bibitem{Steinberg}
R. Steinberg, Invariants of finite reflection groups,
\emph{Canad. J. Math.} 12 (1960), 616--618.
\href{https://doi.org/10.4153/CJM-1960-055-3}{DOI}.

\bibitem{Fausk}
H. Fausk, Equivariant homotopy theory for pro-spectra,
\emph{Geom. Topol.} 12 (2008), 103--176, Appendix A.
\end{thebibliography}
\end{document}
